\section{The finite-chain spectral mesh}\label{app:xx-mesh}
The Hilbert spectral-counting approach goes back to Widom~\cite{Widom66}; see also Fedele~\cite{Fedele20}. Otte~\cite[Sections~3--4]{Otte22} uses the Hilbert-to-Carleman reduction
for determinant asymptotics. The unit-cell comparison and periodization
below give the ordered spectral estimate needed for the finite XX chain.

\begin{lemma}[Finite Hilbert sections and the Carleman kernel]
\label{lem:hilbert-mesh}
For $m\ge1$ and fixed $a>0$ let
\[
H_a^{(m)}=\left(\frac1{j+k+a}\right)_{j,k=0}^{m-1},
\qquad L_m=\ln(2m+1).
\]
Write $\sigma_j(A)$ for the nonincreasing singular values of $A$,
padding finite lists by zeros. Write $(h_j^{(m)})_{j\ge1}$ for the
decreasing rearrangement of
$\{\pi\sech(2\pi^2 k/L_m):k\in\mathbb Z\}$. There is a constant
$C_a$ independent of $m$ with
\[
\sum_{j\ge1}\left|\sigma_j(H_a^{(m)})-h_j^{(m)}\right|\le C_a.
\]
\end{lemma}

\begin{proof}
The difference $H_a^{(m)}-H_2^{(m)}$ is sign definite, since its
entries are the moments of $(t^{a-1}-t)\,dt$ on $[0,1]$. Hence
\[
\|H_a^{(m)}-H_2^{(m)}\|_1
=\sum_{j=0}^{m-1}\left|\frac1{2j+a}-\frac1{2j+2}\right|
\le C_a'.
\]
Consider the integral operator $C_m$ on
$L^2([1/2,m+1/2],dx)$ with kernel $(x+y)^{-1}$. The Laplace
representation
$(x+y)^{-1}=\int_0^\infty e^{-t x}e^{-t y}\,dt$ makes $C_m$ a
positive trace-class operator, with trace
$\int_{1/2}^{m+1/2}(2x)^{-1}dx=L_m/2$.

Compress $C_m$ to the orthonormal indicators of the unit intervals
$[j+1/2,j+3/2)$, $0\le j<m$. The resulting matrix has entries
\[
\int_0^\infty e^{-t(j+k+2)}
\left(\frac{2\sinh(t/2)}{t}\right)^2\,dt.
\]
The factor in parentheses is at least one. The difference from
$H_2^{(m)}$ is positive semidefinite, as is seen by evaluating
quadratic forms inside the Laplace integral. The min--max principle
for compression and positive perturbation therefore gives
$\sigma_j(C_m)\ge\sigma_j(H_2^{(m)})$ for every $j$, with zero
padding. Their difference in $\ell^1$ is
\[
\Tr C_m-\Tr H_2^{(m)}
=\frac12\left(\ln(2m+1)-\sum_{j=1}^m\frac1j\right)=O(1).
\]

The unitary change of variables $x=e^s/2$ transforms $C_m$ to the
convolution kernel $K(s-t)$ on $[0,L_m]$, where
$K(z)=(2\cosh(z/2))^{-1}$. For $L=L_m$, define the periodized
operator $P_L$ by the kernel $\sum_{k\in\mathbb Z}K(s-t+kL)$ on
the same interval. The Fourier transform with convention
$\widehat K(\xi)=\int_{\mathbb R}K(z)e^{-i\xi z}\,dz$ is
$\pi\sech(\pi\xi)$. This follows by setting $u=e^z$ in the
integral and using
$\int_0^\infty u^{b-1}(1+u)^{-1}du=\pi/\sin(\pi b)$ at
$b=1/2-i\xi$. The Fourier modes on $[0,L]$ diagonalize $P_L$,
with eigenvalues $\widehat K(2\pi k/L)$ for $k\in\mathbb Z$.
They are nonnegative and summable for each $L$.

The trace norm of the difference between $P_L$ and the original
convolution operator is bounded uniformly for $L\ge1$. For
$k\ge1$, the image kernel $K(s-t+kL)$ is, after reflection of the
second variable, $K(s+t+(k-1)L)$. Negative images are transposes.
For $k=1$, expand for $z>0$
\[
K(z)=\sum_{j\ge0}(-1)^j e^{-(j+1/2)z}.
\]
Group consecutive terms into rank-two kernels
$f_b\otimes f_b-f_{b+1}\otimes f_{b+1}$ with
$f_b(x)=e^{-b x}$ and $b=2j+1/2$. On $[0,L]$,
$\|f_b\|_2\le (2b)^{-1/2}$ and
$\|f_b-f_{b+1}\|_2\le\|x e^{-b x}\|_{L^2(0,\infty)}=O(b^{-3/2})$.
Each paired trace norm is thus $O(b^{-2})$, uniformly in $L$.
The series converges in trace norm and its kernel agrees almost
everywhere with $K(s+t)$. For $k\ge2$, there is a positive offset
$d=(k-1)L\ge1$. The trace norms of the unpaired rank-one kernels
are bounded in sum by
\[
\sum_{j\ge0}\frac{e^{-(j+1/2)d}}{2j+1}\le C e^{-d/2}.
\]
These bounds sum over the images. The trace-norm perturbation
inequality for singular values now bounds the $\ell^1$ distance
between the spectra of $C_m$ and $P_{L_m}$ by an absolute constant.
Combining with the compression comparison proves the claim.
\end{proof}

\begin{corollary}[Quantitative active-mode mesh of the finite open XX chain]
\label{cor:xx-mesh}
Let $N=2n\ge4$ be even, $m_1=\lceil n/2\rceil$,
$m_2=\lfloor n/2\rfloor$. For the cross block $B_N$ of the
half-filled finite open XX covariance, the singular values are at
bounded $\ell^1$ distance from the decreasing rearrangement of
\begin{equation}
\left\{\frac12\sech(2\pi^2 k/L_{m_i}):
              i=1,2,\ k\in\mathbb Z\right\}.
\label{eq:xx-explicit-mesh}
\end{equation}
For each fixed $0<a<b<\infty$, among the mesh entries whose
energy $4\pi^2|k|/L_{m_i}$ belongs to $[a,b]$, all but $O_{a,b}(1)$
can be matched to finite-chain modes with energies in $[a/2,2b]$.
The sum of the absolute energy errors over those matches is $O_{a,b}(1)$.
\end{corollary}

\begin{proof}
Write $M=N+1$. In reversed left and forward right coordinates
$1\le u,v\le n$ the covariance cross block, up to a global sign
from choosing the other occupied half of the standing waves, has
entries
\[
\frac{\sin(\pi (u+v-1)/2)}{2M\sin(\pi(u+v-1)/(2M))}
-\frac{\sin(\pi(M+v-u)/2)}{2M\cos(\pi(v-u)/(2M))}.
\]
This follows from the product-to-sum identity for the sine waves
and the finite cosine sum. Subtract the principal matrix with
entries $\sin(\pi(u+v-1)/2)/[\pi(u+v-1)]$. The two remaining
reciprocal factors have the form $M^{-1}f(a_u+b_v)$, with
$f(z)=(2\sin(\pi z/2))^{-1}-(\pi z)^{-1}$ and
$|a_u|,|b_v|\le1/2$ for the first term. For the second use
$f(z)=(2\cos(\pi z/2))^{-1}$ with centered coordinates
$|a_u|,|b_v|\le1/4$. The functions are analytic on the disks
of radii $2$ and $1$, respectively, with a removable singularity
at zero for the first. If $|a_u|,|b_v|\le r$, the trace norm
of the degree-$k$ monomial matrix is at most $(n/M)(2r)^k$
by the binomial expansion into rank-one matrices. Each sine is a
sum of two products of phases in $u$ and $v$. Their diagonal-unitary
multipliers preserve these bounds. Absolute summability of the
Taylor coefficients on the relevant smaller radii gives a uniformly
bounded trace-norm error.

Separating the parities and removing the alternating signs makes
the principal matrix equivalent in singular values to the direct
sum of $(2\pi)^{-1}H_{1/2}^{(m_1)}$ and $(2\pi)^{-1}H_{3/2}^{(m_2)}$.
Apply Lemma~\ref{lem:hilbert-mesh} to the
blocks, together with the trace-norm perturbation inequality.
Rearranging the union in decreasing order only decreases the
$\ell^1$ distance, which proves~\eqref{eq:xx-explicit-mesh}.

A true singular value $s$ equals $\sqrt{\lambda(1-\lambda)}$
for an eigenvalue $\lambda$ of the restricted covariance, since
the full covariance is a projection. The one-mode energy
$|\ln(\lambda/(1-\lambda))|$ is therefore
$e(s)=2\arcosh(1/(2s))$, with $e(0)=\infty$.
Match the ordered singular-value lists from~\eqref{eq:xx-explicit-mesh}. Mesh values of
energy in $[a,b]$ have positive distance from the complement
of the singular-value interval corresponding to $[a/2,2b]$.
Thus at most a bounded number of their matches are outside that
interval. On the remaining matches $e(s)$ has bounded derivative,
which converts the $\ell^1$ singular-value bound to the stated
energy bound.
\end{proof}

\begin{corollary}[Fourier damping for tilted energy laws]
\label{cor:xx-fourier}
Let $e_j$ denote the finite open XX mode energies and $W=\ln N$. Under a tilt
$\beta$ in a fixed compact subinterval of $(0,\infty)$, let $X_j$
be independent Bernoulli variables of parameters
$(1+e^{\beta e_j})^{-1}$, with deterministic zero contribution from
any infinite-energy mode. For $\varphi_\beta(t)
=\mathbb E_\beta\exp(it\sum_j e_j X_j)$ there are constants
$c,\delta,\tau>0$ such that for all sufficiently large even $N$,
\[
|\varphi_\beta(t)|\le
\begin{cases}
\exp(-cWt^2), & |t|\le\delta,\\
\exp(-cW), & \delta\le |t|\le \tau W.
\end{cases}
\]
The constants and large-$N$ threshold can be chosen uniformly in
the given compact set of tilts.
\end{corollary}

\begin{proof}
Use just one mesh branch from~\eqref{eq:xx-explicit-mesh} with $k\ge1$ and mesh energies
in a fixed interval $[a,b]\subset(0,\infty)$. The matched true
energies lie in $[a/2,2b]$ with $O(1)$ deletions and total
absolute energy discrepancy $O(1)$. On these modes, the Bernoulli
variances are uniformly bounded below. The squared modulus of
one characteristic factor is
$1-4q(1-q)\sin^2(te_j/2)$, where $q=(1+e^{\beta e_j})^{-1}$.
It follows that $|\varphi_\beta(t)|$ is bounded by
$\exp(-c_0\sum_j\sin^2(te_j/2))$ using the matched modes only.
For $|t|\le\delta$ with $\delta b$ sufficiently small, each
sine square is bounded below by a constant times $t^2$ and the
number of matches is of order $W$.

For $|t|\ge\delta$ the integral
$\int_a^b\sin^2(tx/2)\,dx$ is bounded below by a strictly positive
constant. It is positive on compact subsets away from zero and
tends to $(b-a)/2$ at infinity. The mesh spacing is
$4\pi^2/L_{m_i}$ with $L_{m_i}=W+O(1)$. The corresponding
Riemann sum for the sine squares differs from this integral
divided by the mesh spacing by $O(1+|t|)$, since the derivative
in $x$ is bounded by $|t|/2$. Deleting $O(1)$ terms and replacing
mesh energies by matched true energies changes the sum by at
most another $O(1+|t|)$ by the same derivative bound. The resulting
lower bound is $c_1W-C_1(1+|t|)$ with $c_1,C_1>0$ independent of $N$.
Choose $\tau\le c_1/(4C_1)$ and take $N$ large enough that
$C_1\le c_1W/4$. The sum is then at least $c_1W/2$ throughout
$\delta\le |t|\le\tau W$, proving the second estimate.
\end{proof}

\begin{corollary}[Transfer to the finite open critical Ising chain]
\label{cor:ising-fourier}
Let $L$ be even. Across its middle cut, the entanglement-mode
energies of the finite open critical Ising chain on $L$ spins
occur twice in the finite open XX chain on $2L$ spins across
its middle cut. In particular, the XX bounds in
Corollary~\ref{cor:xx-fourier} apply to the Ising characteristic
function with adjusted constants.
\end{corollary}

\begin{proof}
This is the covariance decomposition behind the even-block
identity of Igl\'oi and Juh\'asz~\cite[eqs.~(19)--(23)]{IgloiJuhasz08}.
Write $a_j=c_j+c_j^\dagger$ and $b_j=i(c_j^\dagger-c_j)$
for the two Majoranas of the XX site $j$. Its hopping term is
$\tfrac i2(a_j b_{j+1}-b_j a_{j+1})$. There are two disjoint
Majorana paths, one using $a_j$ on odd sites and $b_j$ on even
sites and the other using the opposite assignment. In each path,
the signs of the hopping terms can be made uniform by changing
signs of individual Majoranas. Each path is then, up to a positive
scale of its Hamiltonian, the critical Ising chain of $L$ spins
in its Majorana representation. Since $L$ is even, the first
$L$ XX sites restrict each path to its first $L$ Majoranas, which
are precisely the first $L/2$ Ising spins. The restricted Gaussian
covariance is the direct sum of the two path covariances. The
mode-energy lists therefore add by multiset union. The tilted
characteristic function of the XX energy is the square of the
Ising characteristic function. Taking its modulus square root
gives the desired damping.
\end{proof}
